% !TEX root = ../../main.tex

\section{Objetivos}

El objetivo general de esta memoria es ampliar experimentalmente el espacio de planes de ejecución que \texttt{ApplicativeDo} puede evaluar en \textsf{GHC}, considerando reordenamientos válidos de statements dentro de bloques \texttt{do} marcados explícitamente bajo una hipótesis de conmutatividad. Con ello, se busca seleccionar, según el modelo de costos adoptado, un plan de menor costo cuando alguno de estos reordenamientos válidos lo permita.

Los objetivos específicos son los siguientes:

\begin{enumerate}
    \item Caracterizar los reordenamientos válidos de los statements de un bloque \texttt{do} según las dependencias locales que deben preservarse.
  \captionsetup{skip=14pt}
    \item Ampliar el conjunto de planes de ejecución que \texttt{ApplicativeDo} puede considerar para bloques \texttt{do} bajo una hipótesis explícita de conmutatividad, de modo que se pueda seleccionar un plan de ejecución con un costo estructural mínimo entre los obtenidos a partir de reordenamientos válidos.
    \item Comprobar experimentalmente, en los casos conmutativos evaluados, que los reordenamientos válidos preservan el resultado observable del programa y que la selección realizada corresponde a un plan de ejecución con costo estructural mínimo.
\end{enumerate}
